%---------------------------------------------------------
%	7. LIMITATIONS, CONCLUSION & Q&A
%---------------------------------------------------------
\section{Conclusion}

\begin{frame}
    \frametitle{Limitations \& Future Work}
    \begin{columns}[t]
        \begin{column}{0.5\textwidth}
            \footnotesize
            \textbf{Limitations}
            \begin{itemize}
                \item Depthwise convs: few MACs but not GPU-friendly, so speed gains are smaller than compute gains
                \item No runtime benchmark yet (only MACs / parameters)
                \item Not best on every metric (SSIM/LPIPS on LOLBlur, real photos)
                \item Real/synthetic gap in the data (EDA)
                \item Our app: TensorRT disabled for RetinexFormer / HVI-CIDNet (wrong outputs); NightClear-s weights not bundled
            \end{itemize}
        \end{column}
        \begin{column}{0.5\textwidth}
            \footnotesize
            \textbf{Future work}
            \begin{itemize}
                \item Measure latency on phone / edge GPUs; INT8 quantisation
                \item Video: temporal consistency between frames
                \item Collect real paired night data (e.g.\ with a tripod) to close the domain gap
                \item Train jointly with a detector for night-time object detection
            \end{itemize}
        \end{column}
    \end{columns}
\end{frame}

\begin{frame}
    \frametitle{Conclusion}
    \textbf{Key takeaways}
    \begin{enumerate}
        \item \textbf{One model, three problems}: NightClear brightens, denoises and deblurs at once and beats two-model pipelines by +8.7\,dB.
        \item \textbf{Efficient by design}: Fourier-amplitude encoder + dilated decoder give the best LOLBlur PSNR at 7.25 GMACs; our EDA and amplitude-swap test show why.
        \item \textbf{Deployable}: a web app with 5 models, video/webcam and FP16/TensorRT, runnable on CPU or a free GPU.
    \end{enumerate}
\end{frame}
